% Independent reconstruction of GT C.2.13--16 and C.2.18.
\subsection{Selected later failures and the limits of fixed finite tests}

\begin{result}{E.35}{Proved}{An explicit augmentation bound for a five-coordinate group}
\label{coh:augmentation-bound}
Let $p\ge5$ be prime, $a\ge1$, and $M=p^a$.  The universal
two-generator group of nilpotency class at most three and exponent
dividing $M$, denoted $P_M$, has order $M^5$.  If $I$ is the
augmentation ideal of $\mathbb F_p[P_M]$, then
\[
                         I^{1+10(M-1)}=0.
 \tag{E.35.1}
\]
This confirms the exact bound in GT\_TOWER\_CHECK C.2.14.
\end{result}
\begin{proof}
Use the commutator convention $[u,v]=uvu^{-1}v^{-1}$ and put
$c=[x,y]$, $d=[c,x]$, $e=[c,y]$.  The integral Hall model of class
three has five coordinates, with unique normal forms
$x^iy^jc^hd^ue^v$.  Its polynomial denominators involve only two and
three.  Reduction modulo $M$ therefore gives a group with $M^5$
elements.  Equivalently, the truncated Baker--Campbell--Hausdorff
law on the free class-three Lie algebra over $\mathbb Z/M$ gives
that group, with coordinates in weights $1,1,2,3,3$.  Powers are
integer multiples of logarithms, so its exponent divides $M$.
In any class-three exponent-$M$ group, the five Hall generators have
orders dividing $M$; hence every word whose Hall exponents are
multiples of $M$ maps to the identity.  This proves the asserted
universal property as well as the order.

In the group algebra write
\[
 X=x-1,\quad Y=y-1,\quad C=c-1,\quad D=d-1,\quad E=e-1.
\]
The ordered monomials
\[
 X^iY^jC^hD^uE^v,\qquad 0\le i,j,h,u,v<M,
 \tag{E.35.2}
\]
are a basis.  Indeed their binomial expansions in the ordered group
basis are triangular with diagonal one.  Each of $X,Y,C,D,E$ has
$M$-th power zero, since $(g-1)^M=g^M-1$ in characteristic $p$.
The last two variables are central.  The group identities
$yx=xyc^{-1}d^{-1}e^{-1}$, $cx=xcd$, $cy=yce$ give exactly
\[
\begin{aligned}
 YX&=XY+(1+X)(1+Y)
       \big((1+C)^{-1}(1+D)^{-1}(1+E)^{-1}-1\big),\\
 CX&=XC+(1+X)(1+C)D,\\
 CY&=YC+(1+Y)(1+C)E.
\end{aligned}
 \tag{E.35.3}
\]
Every inverse is the finite geometric series in its nilpotent
variable.  Assign weights $1,1,2,3,3$ to $X,Y,C,D,E$.  Every
summand on the right of each rule has weight at least that of the
left.

We justify that this local observation gives a filtration for
arbitrary products.  Move $D,E$ to the right.  First count inversions
of $X,Y$, ignoring occurrences of $C$, and then count occurrences
of a $C$ before an $X$ or $Y$.  Move intervening $C$'s right by the
last two rules when needed, and apply the first rule to a $YX$
inversion.  Every resulting summand has fewer $X,Y$ inversions:
it has either the correctly ordered pair, one letter of the pair,
or neither, and introduces no new $X$ or $Y$.  The last two rules
do not increase the first count and strictly decrease the second.
A first-rule step can create only finitely many $C$'s.  Lexicographic
induction on the two counts therefore terminates in ordered
monomials, throughout without lowering weight.  Terms with an
exponent at least $M$ vanish.

Let $J_n$ be the span of basis monomials of weight at least $n$.
The collection just proved gives $J_rJ_s\subset J_{r+s}$.
The augmentation ideal is $J_1$, since these are precisely all
nonconstant basis monomials.  The largest possible weight in
(E.35.2) is $10(M-1)$, proving (E.35.1).
\end{proof}

\begin{result}{E.36}{Proved}{A finite integral class already detects the incompatible local data}
\label{coh:local-finite-exclusion}
Let $p\ge5$, let $q\ne p$ be a rational prime with $q\equiv1\pmod p$,
and let $m\ne0$ be an integer.  Set
\[
 k=v_p(m),\qquad B(p,m)=10(p^{k+1}-1).
 \tag{E.36.1}
\]
In the tangentially marked geometrically pro-$p$ quotient over
$\mathbb Q_q$ of integral nilpotency class $B(p,m)$, there is no
continuous section having tame inertia with abelian coordinates
$(m,0)$ and geometric Frobenius coordinate a nonzero central element
in the class-two exponent-$p$ quotient $H_p$.
This confirms GT\_TOWER\_CHECK C.2.15 with its stated constant.
\end{result}
\begin{proof}
Put $s=B(p,m)$, $M=p^{k+1}$ and $P=P_M$ from E.35.  Its universal
property makes the quotient from the free pro-$p$ group
characteristic.  The actual arithmetic action therefore descends
to $P$, and reduction of its first three Hall coordinates gives
the actual quotient $P\to H_p$.

Suppose the asserted section exists.  By E.32 the tangential
action at $q$ is unramified at every integral class.  Wild and
non-$p$ tame inertia have zero geometric image.  For tame
$p$-inertia $T$ and arithmetic Frobenius $F$, let $V,U$ be the
section's geometric coordinates.  If $\varphi_q$ is the actual
Frobenius action, the identities are
\[
 U\varphi_q(V)U^{-1}=V^q,\qquad
 \varphi_q(x)=x^q,\qquad
 \operatorname{gr}_i(\varphi_q)=q^i.
 \tag{E.36.2}
\]
Work temporarily in the rational class-$s$ Hall group over
$\mathbb Q_p$.  Successive central corrections, dividing in weight
$i$ by the nonzero scalar $q^i-1$, give $W$ with
$W^{-1}U\varphi_q(W)=1$.  Thus $V'=W^{-1}VW$ satisfies
$\varphi_q(V')=(V')^q$ and has abelian image $(m,0)$.
Compare $V'$ with $x^m$.  At the first differing central layer of
weight $i\ge2$, their difference is killed by $q^i-q$, which is
nonzero in $\mathbb Q_p$.  Induction gives
\[
                         W^{-1}VW=x^m.
 \tag{E.36.3}
\]

Only the integral consequences of this conjugacy will be evaluated
in a finite group algebra.  The truncated Magnus series of $V$
and $x^m$ have coefficients in $\mathbb Z_p$.  Rational conjugacy
implies equality of the sums of their coefficients on each cyclic
orbit of words through degree $s$.  The integral quotient of the
word module by cyclic rotations is free, with those orbits as
basis.  Therefore the integral difference of the two series is
an integral sum of additive algebra commutators, through degree
$s$.  This is the finite-level form of the argument in A.30.

Evaluate those integral expressions, and only those expressions,
at $X=x-1,Y=y-1$ in $\mathbb F_p[P]$.  E.35 gives
$I^{s+1}=0$, so higher-degree terms vanish.  The images of $V$
and $x^m$ are consequently equal in
$\mathbb F_p[P]/[\mathbb F_p[P],\mathbb F_p[P]]$.
This quotient has the conjugacy classes of $P$ as basis: an
additive commutator of group elements is a difference within one
conjugacy class, and every conjugate difference has this form.
Equality of two group basis elements implies that they are
conjugate.  We obtain $h\in P$ with $h^{-1}Vh=x^m$.
There has been no attempt to evaluate the rational element $W$
in characteristic $p$.

After geometric conjugation by $h$ the Frobenius coordinate is
$U'=h^{-1}U\varphi_q(h)$.  On $H_p$ the actual class-two action
is the cyclotomic action and is the identity since $q\equiv1$
modulo $p$.  The prescribed central image of $U$ is therefore
unchanged.  Equation (E.36.2) in $P$ now gives
\[
                         U'x^{qm}(U')^{-1}=x^{qm}.
 \tag{E.36.4}
\]
Use the class-three Lie basis over $\mathbb Z/M$
\[
 X_0,\ Y_0,\ C_0=[X_0,Y_0],\
 D_0=[C_0,X_0],\ E_0=[C_0,Y_0],
\]
and write
$\log U'=aX_0+bY_0+c_0C_0+d_0D_0+e_0E_0$.
The finite exponential of its adjoint action yields
\[
 \log(U'x^{qm}(U')^{-1})
 =qmX_0-qmbC_0+
 qm(c_0+ab/2)D_0+qmb^2E_0/2.
 \tag{E.36.5}
\]
Since $v_p(qm)=k$, the $C_0$ coordinate of (E.36.4) forces
$b\equiv0\pmod p$, and the $D_0$ coordinate then forces
$c_0\equiv0\pmod p$.  In $H_p$, where a central element also
has zero $a,b$ coordinates, this rules out a nonzero central
image of $U'$.  That contradicts the preserved Frobenius image,
and proves the claim.
\end{proof}

We will use the following actual finite specialization, without
requiring equations for a generic Heisenberg cover.  For every
$p\ge5$, reduce the point-three section in the tangential class-two
exponent-$p$ quotient.  It gives a continuous homomorphism
\[
 \rho_{3,p}:G_{\mathbb Q}\longrightarrow
                 H_p\rtimes\mathbb F_p^\times.
 \tag{E.37.0}
\]
Its image is the entire group and the maximal abelian subfield
of its fixed field $L_p$ is $\mathbb Q(\mu_p)$.  Indeed, over
$\mathbb Q(\mu_p)$ the Kummer characters of $3,-2$ are independent:
their valuations at three and two are one, and these primes are
unramified in that cyclotomic field.  The image therefore surjects
onto $H_p^{\rm ab}$.  A subgroup with that property is $H_p$
itself, because commutators of two lifts with independent abelian
coordinates generate its center.  The cyclotomic quotient is
surjective as well.  Finally, the derived subgroup of
$H_p\rtimes\mathbb F_p^\times$ is exactly $H_p$: conjugation by
the scalar two produces all its abelian coordinates through the
invertible scalar $2-1$, and their commutators produce the center.
This proves the maximal-abelian assertion with the actual marking.

Also, for every $p\ge5$,
\[
 H^1(G_{\mathbb Q,\{p,\infty\}},\mathbb F_p(2))=0.
 \tag{E.37.1}
\]
Inflation injects the integral first cohomology into the absolute
group, which is zero at $p$ by A.26.  Its integral second cohomology
is $K_2(\mathbb Z)\{p\}=0$ by A.22.  The multiplication-by-$p$
continuous coefficient sequence gives (E.37.1).

\begin{result}{E.37}{Proved}{Selected fixed-depth examples have a bounded later first failure}
\label{coh:selected-fixed-depth-failure}
For every $R\ge1$ and $N\ge1$, the construction in E.33 can be
chosen to give an ordinary infinite-support pair that survives
full integral class $R$ with the specified finite Magnus
congruence, but whose first empty unrestricted integral section
fibre is a finite class $r>R$.  More precisely, let
\[
 R_0=\max(R,3),\quad d=R_0-1,\quad n=J!,\quad
 a=v_5(n),
\]
and choose the permitted constant $C$ additionally divisible by
five.  At a distinguished first positive support prime $q$, put
$m=E_0(q)$.  Then
\[
 v_5(m)=v_5(C)+da+v_5(d!),\qquad
 R_0<r\le10(5^{v_5(m)+1}-1).
 \tag{E.37.2}
\]
The pair has no full geometrically pro-five section.  The
point-three twist of the characteristic Heisenberg cover of
geometric degree $125$, fixed independently of $R,N$, detects
that full pro-five exclusion.  This confirms GT\_TOWER\_CHECK
C.2.13 and C.2.16.  It determines a two-sided bound, not the
exact least class within that interval.
\end{result}
\begin{proof}
Run E.33 for $R_0,N$, enlarging $C$ by five.  At the first arithmetic
stage of E.31 let $P$ be the complete CRT modulus, and let
$\sigma_F$ be the prescribed automorphism of
$F=\mathbb Q(\mu_P)$.  It fixes $\mu_n$, since its CRT class is
one modulo $n$.  Here $5\mid n$.  We impose one additional
condition on the first positive support prime.

Let $L_5$ be the actual specialization field from (E.37.0), and put
\[
 R_b=F((-2)^{1/n}),\qquad R'=R_b(3^{1/5}).
 \tag{E.37.3}
\]
E.31 gives an extension of $\sigma_F$ to $R_b$ fixing the displayed
$n$-th root: the original radical field over $\mathbb Q(\mu_n)$
has that field as its maximal abelian subfield.  The fifth-power
classes of $3,-2$ are independent in any cyclotomic field $F$.
For a nontrivial rational combination its valuation at two or
three is not divisible by five.  Its real fifth root has degree
five and generates a nonnormal real field, which cannot be a
subfield of an abelian extension of $\mathbb Q$.
The cyclic extension $R_b/F$ has its unique degree-five
subextension $F((-2)^{1/5})$.  Thus
$R_b\cap F(3^{1/5})=F$, and $\sigma_F$ extends to an
automorphism $\sigma_{R'}$ fixing both displayed radicals.

Since $L_5\cap F=\mathbb Q(\mu_5)$, we have
$\operatorname{Gal}(L_5F/F)=H_5$.  The abelian extension
$R'/F$ contains its full maximal abelian subextension.
Therefore, exactly,
\[
                 L_5F\cap R'=F(3^{1/5},(-2)^{1/5}).
 \tag{E.37.4}
\]
Choose any nonzero central element $z$ of $H_5$.  Its action on
$L_5$ and $\sigma_F$ on $F$ agree on their intersection.
The combined action on $L_5F$ agrees with $\sigma_{R'}$ on
(E.37.4), since both fix the two radicals and have the same
restriction to $F$.  These are compatible automorphisms of the
actual compositum.  Chebotarev therefore supplies arbitrarily
large first primes $q$ satisfying every original CRT and radical
condition, also splitting $3^{1/5}$, and with
$\rho_{3,5}(\operatorname{Frob}_q)$ central nonzero.
All stated conditions are conjugacy invariant; conjugation may
scale $z$ by a nonzero scalar and cannot make it zero.
Choose such a $q$ and continue the arithmetic recursion unchanged.
E.31 and E.33 retain all their conclusions for this same pair.

We show that every integral class-two section above it has the
same $H_5$ reduction as the point-three section.  Choose the
literal convergent Kummer cocycles of the constructed products.
Every fresh exponent is divisible by five, so their reductions
are the literal cocycles of $(3,-2)$.  After aligning abelian
cocycles, the difference of the central coordinates is a
$\mathbb F_5(2)$-cocycle.  At a fresh support prime its inertia
is zero: E.32 gives central inertia in $(E_q/2)\mathbb Z_5$,
and the point-three inertia is zero.  At any other prime greater
than three and different from five, both integral inertia
coordinates are zero by the tame weight relation.  At two and
three, the difference is killed on inertia by $q'(q'-1)$,
a unit modulo five.  The difference is consequently unramified
outside five and is a coboundary by (E.37.1).

At the selected $q$, the cyclotomic action modulo five is the
identity, so a central coboundary vanishes on its Frobenius.
Its geometric Frobenius image is therefore the same nonzero
central element for every integral lift.  Geometric conjugation
also preserves that value.  Its inertia abelianization is $(m,0)$.
E.36 excludes a section at integral class
$10(5^{v_5(m)+1}-1)$ already after restriction to
$G_{\mathbb Q_q}$ and projection to the pro-five quotient.
Hence the full integral fibre at that class is empty.
It is nonempty through class $R_0$ by E.33; projection makes
nonemptiness downward closed.  A least empty class therefore
exists in the interval (E.37.2), and no full pro-five section can
exist.

The exponent is exact.  The first-stage shell in E.31 gives
$v_5(q-1)=v_5(n)=a\ge1$.  For $i\ge1$,
$v_5(q^i-1)=a+v_5(i)$: the assertion for $i$ prime to five follows
by summing a geometric series, and passage from $i$ to $5i$
adds one by the binomial expansion of a number congruent to one
modulo five.  Applying this to
$m=Cn\prod_{i=2}^d(q^i-1)$ proves the formula in (E.37.2).

For completeness we give the promised cover comparison in the
same actual marking.  Let $K$ be the characteristic kernel of
$\Pi\to H_5$.  The inverse image of the point-three section in
$H_5\rtimes G_{\mathbb Q}$ defines a degree-$125$ cover $Y_5$
with a rational point above three.  Its geometric pro-five
topology is the full induced topology, by the prime-power-index
comparison C.3.  The preceding class-two calculation would
geometrically conjugate any hypothetical full pro-five section
above our pair into this cover subgroup.

At the selected prime, the ordinary local point $t_q$ has positive
valuation $m$.  It is a fifth power in $\mathbb Q_q$: its class
relative to three is a fifth power, and $q$ was chosen to split
$3^{1/5}$.  Also $1-t_q$ is a principal $q$-unit and hence a
fifth power.  The actual point section thus has zero abelian
Frobenius coordinates modulo five, and its central inertia is
$mu$, zero modulo five, for an integral local primitive $u$.
Its Frobenius image in $H_5$ is central.  It cannot be nonzero,
because the point section exists in every integral class and
would contradict E.36.  Therefore its entire local $H_5$ image
is the identity.  This argument establishes the needed local
marking comparison without an unproved generic cover equation.

The local $Y_5$-fibre is the coset set $H_5$.  The identity image
of the ordinary point section and the nonzero central
point-three Frobenius make Frobenius act there by a nontrivial
central translation, which has no fixed coset.  A full pro-five
section with the given nonzero ordinary valuation would be
locally conjugate to that point section by A.31, and its global
factorization through $Y_5$ would supply a fixed coset, a
contradiction.  This gives the stated single fixed-cover
detection as well as the earlier explicit finite-class bound.
\end{proof}

\begin{result}{E.38}{Proved}{Any prescribed finite geometric test can coexist with first failure at class three}
\label{coh:arbitrary-finite-first-three}
Let $N\triangleleft\Pi$ be open and normal in the full arithmetic
group, and put $Q=\Pi/N$.  There exists an ordinary
infinite-support completed pair with a continuous section of
\[
 E_Q\times_{E_{Q^{\rm ab}}}E_{\rm ab}
       =E_U/(N\cap\overline{[\Pi,\Pi]}),
 \tag{E.38.1}
\]
whose $Q$-component is the point-three section, and whose
restriction at every finite prime is geometrically conjugate
to the image of its ordinary point; at the real place it is
the image of three.  Nevertheless its unrestricted full
integral class-two fibre is nonempty and its class-three
fibre is empty.

Thus, for every one fixed finite geometric quotient, existence
with the joint full abelianization and all those local
comparisons does not imply integral class-three liftability.
This confirms the exact quantifiers of GT\_TOWER\_CHECK C.2.18.
\end{result}
\begin{proof}
Use the fully reconstructed finite-cover recursion E.20, and
retain its point-three cover $Y$, finite Galois boundary field
$F$, degree $M=[F:\mathbb Q]$, bad-prime integer $\Delta$, and
the conductors $f_3,f_{-2}$ of its fixed auxiliary radical fields.
These data are fixed before the following new coefficient prime.
Put
\[
 B_0=\operatorname{lcm}(24,4M,|Q|,\Delta,f_3,f_{-2}),
 \quad p\ge7,\quad p\nmid B_0,\quad N_0=pB_0,
 \tag{E.38.2}
\]
and take $C=\operatorname{lcm}(24,4M,|Q|,p)$.  These are allowed
parameters of E.20, since $C\mid N_0$.  Its exponents are
$n_j=N_0(j+4)!$, and $v_p(n_0)=1$.

At the initial stage choose the actual rational point over three
with Kummer coordinates $(1,1)$ at every old local place.
The first CRT class is therefore one at every prime.  If $P$ is
its modulus, E.20 permits the usual required precisions, and the
precision at $p$ is at least
$2v_p(N_0\,6!)+4=6$.  In particular
\[
                    q\equiv1\pmod P,\qquad p^2\mid P
 \tag{E.38.3}
\]
is the original allowed condition on the first positive support
prime.  We add one compatible central Frobenius condition.

Let $L_p$ be the actual point-three field (E.37.0), and write
\[
 K=F\mathbb Q(\mu_P),\qquad
 R_b=K((-2)^{1/n_0}),\qquad R'=R_b(3^{1/p}).
 \tag{E.38.4}
\]
The maximal-abelian assertion after (E.37.0) gives
$L_p\cap\mathbb Q(\mu_P)=\mathbb Q(\mu_p)$, and
$L_p\mathbb Q(\mu_P)/\mathbb Q(\mu_P)$ has group $H_p$.
The Galois extension $K/\mathbb Q(\mu_P)$ has degree dividing
$M$, prime to $p$.  The two extensions are linearly disjoint.
Consequently
\[
 \operatorname{Gal}(L_pK/K)=H_p,\qquad
                  L_p\cap K=\mathbb Q(\mu_p).
 \tag{E.38.5}
\]

The classes of $3,-2$ are independent modulo $p$-th powers in
$\mathbb Q(\mu_P)$, by the nonnormal-real-radical argument used
in E.37, now with degree $p$.  They remain independent in $K$:
restriction followed by norm is multiplication by the
prime-to-$p$ extension degree.  The cyclic extension $R_b/K$
therefore has unique degree-$p$ subextension
$K((-2)^{1/p})$, and
$R_b\cap K(3^{1/p})=K$.  The abelian extension $R'/K$ contains
the full maximal abelian subextension of $L_pK/K$, so
\[
                   L_pK\cap R'=K(3^{1/p},(-2)^{1/p}).
 \tag{E.38.6}
\]
Take a nonzero central element on $L_p$, the identity on $K$,
and the identity on $R'$.  Equations (E.38.5)--(E.38.6) prove
that these restrictions agree on every intersection and define
an automorphism of the actual compositum.  Chebotarev supplies
arbitrarily large primes $q$ in its conjugacy class.  Such a
prime satisfies (E.38.3), splits in $F$ and in the original
first radical field
$\mathbb Q(\mu_{n_0},(-2)^{1/n_0})$, and its point-three
$H_p$ Frobenius is central nonzero.  Conjugation only rescales
that nonzero central value.  Thus it is a legal first prime of
E.20 with the added condition.  Continue every remaining choice
by E.20.  Its component maps and compact local conjugator argument
give exactly the section and local comparisons in (E.38.1)
for the ordinary pair
\[
             A=3\prod_{j\ge0}q_j^{n_j},\qquad
             B=-2\prod_{j\ge0}r_j^{n_j}.
 \tag{E.38.7}
\]

Every $n_j$ is divisible by $p$.  After aligning literal abelian
cocycles, the reduction of any integral class-two section differs
centrally from point three by a $\mathbb F_p(2)$-cocycle.
At a new support prime its central inertia is the valuation times
an integral local primitive, hence zero modulo $p$.  At all other
primes greater than three and different from $p$, both inertia
values are integrally zero.  At two and three the difference is
killed by $q'(q'-1)$, a unit modulo $p$.  It is therefore
unramified outside $p$ and is a coboundary by (E.37.1).
Every integral class-two lift has the point-three $H_p$ image.

At our selected $q$, its actual Frobenius coordinates consequently
satisfy
$\alpha\equiv\beta\equiv0\pmod p$ and
$\gamma\not\equiv0\pmod p$.  Geometric conjugation preserves these
properties because $q\equiv1\pmod p$ and this image is central.
The ordinary valuation vector is $(m,0)$ with $m=n_0$ and
$v_p(m)=1$.  Locally the second Kummer cocycle is an integral
coboundary $b=\delta u$; its primitive is integral because
$1-t_q$ is a principal $q$-unit.  The exact class-two tame
identity gives $\beta=(q-1)u$ and $h(T)=mu$ for every integral
lift.  B.23 therefore gives its actual first weight-three residue
\[
 Z_{D,q}=m(\gamma-q\alpha u)
                  +(q^2-1)u\binom m2.
 \tag{E.38.8}
\]
The first summand has valuation exactly one, while the second has
valuation at least three: $p^2\mid(q-1)$, $p\ge7$, and
$v_p\binom m2=1$.  Hence $v_p(Z_{D,q})=1$ but
$v_p(q^2-1)\ge2$.  This is a nonzero local integral obstruction,
already nonzero modulo $p^2$.  No integral class-two section
above the pair can lift to class three.  A.25 supplies integral
class-two existence, so the first unrestricted integral failure
is exactly class three.

The last argument concerns reductions of every integral
class-two section.  It does not classify arbitrary finite
coefficient class-two sections, which may have additional
ramification and central classes.  The locally matching section
of (E.38.1), on the other hand, is the actual one proved by
E.20, with the full abelianization retained; it is not merely a
section modulo a further finite coefficient.  These are the
distinct quantifiers needed for the assertion.
\end{proof}
